\chapter{Calcul différentiel et optimisation}
\label{ch:jour2}

\begin{aitip}
Ce chapitre est en cours de rédaction. Le cours est esquissé ci-dessous ; le texte complet sera publié dans la prochaine itération de l'atelier. Le syllabus et le notebook de TP de cette journée sont déjà disponibles.
\end{aitip}

\section*{Sections prévues}
\begin{itemize}
  \item Dérivées multivariées et règle de la chaîne.
  \item Descente de gradient, descente de gradient stochastique, Adam.
  \item Rétropropagation comme règle de la chaîne sur un graphe de calcul.
  \item Programmation différentiable : pourquoi un framework moderne calcule des gradients que vous n'avez pas écrits.
  \item Taux d'apprentissage, hyperparamètres et lecture d'une courbe d'apprentissage.
\end{itemize}

\section*{Travail pratique}
Le notebook du Jour 2 (\texttt{Day2\_Calculus\_Optimization.ipynb}) parcourt l'implémentation de la rétropropagation à la main en NumPy pour le MLP du Jour 1, puis entraîne un réseau PyTorch plus profond de bout en bout et inspecte les courbes d'apprentissage qui en résultent.
